\documentclass[11pt]{article}
\usepackage[letterpaper,margin=0.9in]{geometry}
\usepackage[hidelinks]{hyperref}
\usepackage{enumitem}
\usepackage{titlesec}
\usepackage{tabularx}

\setlist[itemize]{leftmargin=1.2em,itemsep=3pt,topsep=3pt}
\titleformat{\section}{\large\bfseries}{}{0em}{}[\titlerule]

\newcommand{\entry}[4]{
  \noindent\begin{tabularx}{\textwidth}{@{}Xr@{}}
    \textbf{#1} & #2 \\
    \textit{#3} & \textit{#4}
  \end{tabularx}\vspace{2pt}
}

\begin{document}

\begin{center}
  {\LARGE \textbf{Piter Z. Garcia Bautista}}\\[4pt]
  Rochester, NY \textbar\ +1 (646) 288-3300 \textbar\ \href{mailto:pzg8794@rit.edu}{pzg8794@rit.edu}\\
  \href{https://linkedin.com/in/garciapiterz}{linkedin.com/in/garciapiterz} \textbar\
  \href{https://github.com/pzg8794}{github.com/pzg8794} \textbar\
  \href{https://garciapiterz.com}{garciapiterz.com}
\end{center}

\section{Research Profile}
Graduate researcher focused on deep learning and computational biology, with hands-on experience in neural network-based approaches to biological sequence and structure data. Background includes RNA secondary structure prediction with thermodynamic deep learning enhancements (BIO614), RNA-seq differential-expression analysis (BIO550), fairness-aware ML applied to bioinformatics diagnostics (ISTE-780), and reproducible evaluation pipeline design for biological data. Current Graduate Assistant work (Quantum and AI Research, RIT) deepens expertise in neural architectures, uncertainty quantification, and structured experimental methodology. Motivated to work on deep learning methods that unlock biological insights from complex, high-dimensional, and noisy molecular data.

\section{Research Interests}
Deep learning for biological sequence and structure analysis; neural network methods for genomics and transcriptomics; computational approaches to RNA structure prediction and differential expression; fairness-aware and reproducible ML for biomedical data; uncertainty quantification in biological deep learning models.

\section{Education}

\entry{Rochester Institute of Technology}{Expected Aug 2026}{Master of Science, Data Science --- GPA 3.9+}{Rochester, NY}
\begin{itemize}
  \item Concentration: Bioinformatics \& Artificial Intelligence.
  \item Relevant coursework: Bioinformatics Algorithms (BIOL-630), Neural Networks \& Deep Learning, Data-Driven Knowledge Discovery (ISTE-780), Applied Data Science (DSCI-601), Foundations of Data Science (DSCI-633), Software Engineering for Data Science (DSCI-644).
\end{itemize}

\entry{University of Rochester, Warner School of Education}{In progress}{Graduate study, Teaching Computer Science K--12 (M.S. track)}{Rochester, NY}
\begin{itemize}
  \item NSF Robert Noyce Scholar (\$30,000 full funding); Advanced Certificate in Leadership in Disability and Inclusive Practices.
\end{itemize}

\entry{Rochester Institute of Technology}{2015}{Master of Science, Computer Science}{Rochester, NY}
\begin{itemize}
  \item Concentrations: Artificial Intelligence, Big Data Analytics, Computer Graphics.
  \item MS Thesis: ASL Recognition Application --- image processing, feature extraction, and neural pattern classification.
\end{itemize}

\entry{Farmingdale State College}{2012}{Bachelor of Science, Computer Engineering}{Farmingdale, NY}
\begin{itemize}
  \item Dual degree: Electrical Engineering \& Computer Engineering Technology.
  \item Honors: CSTEP Award \& Merit; Dual BS Scholarship.
\end{itemize}

\section{Research and Professional Experience}

\entry{Rochester Institute of Technology, Software Engineering Department}{Fall 2025 -- Present}{Graduate Assistant, Quantum and AI Research}{Rochester, NY}
\begin{itemize}
  \item Build reproducible ML and deep learning evaluation pipelines across local, Colab, and GCP environments, with structured logging and rigorous benchmarking methodology.
  \item Develop and evaluate neural network architectures for sequential and structured data, with emphasis on reliability, reproducibility, and generalization under distribution shift.
  \item Produce manuscript-ready technical analyses of deep learning methods including evaluation under noisy, heterogeneous, and partially observed data conditions.
\end{itemize}

\entry{Rochester Institute of Technology}{2025 -- 2026}{RNA Structure Prediction and Biological Deep Learning (BIO614 and BIO550)}{Rochester, NY}
\begin{itemize}
  \item Led substantial computational work in BIO614: advanced RNA secondary structure prediction by integrating the Nussinov dynamic-programming algorithm with thermodynamic MFE enhancements using Turner parameters and ViennaRNA, achieving $>$90\% accuracy on synthetic motifs.
  \item Extended prediction pipeline with LSTM and neural-network-based scoring components to capture long-range dependencies in RNA sequences beyond classical dynamic-programming scope.
  \item Produced formal manuscript-style project (\textit{BIO614-FinalProjectProposal}, Overleaf) with reproducible implementation, structured evaluation, and biological interpretation.
  \item In BIO550: conducted RNA-seq differential-expression analysis, including reproducible preprocessing, quality-control validation (DESeq2 workflow), and biological interpretation of gene-expression patterns.
\end{itemize}

\entry{Rochester Institute of Technology}{2024 -- 2025}{Equitable Bioinformatics Research --- ISTE-780}{Rochester, NY}
\begin{itemize}
  \item Pioneered fairness assessment applied to RNA secondary structure prediction algorithms, implementing bias-detection frameworks using SHAP analysis and disparate-impact metrics.
  \item Developed 900+ line modular Python codebase for reproducible fairness-aware evaluation of bioinformatics ML methods.
  \item Demonstrated statistical significance ($p < 0.01$) in algorithmic bias detection across demographic proxies in biological data.
\end{itemize}

\entry{VEDADATA Inc.}{Sep 2022 -- Aug 2024}{Data Solutions Engineer}{New York, NY}
\begin{itemize}
  \item Designed scalable Python data pipelines for ingestion, validation, statistical diagnostics, and analytics operations on large biological and health datasets.
\end{itemize}

\entry{VIOME Inc.}{Jan 2021 -- Sep 2022}{AI Data Solutions Engineer}{New York, NY}
\begin{itemize}
  \item Developed healthcare AI workflows including deep learning feature engineering, model experimentation, and evaluation for personalized health prediction in AWS environments.
  \item Built preprocessing pipelines for microbiome sequencing data, including normalization, QC filtering, and feature extraction for downstream ML models.
\end{itemize}

\entry{Samasta Health Foundation}{Mar 2017 -- Jan 2021}{Computer \& Data Science Consultant (pro bono)}{Remote}
\begin{itemize}
  \item Designed AI and ML solutions for healthcare equity with a focus on accessible, evidence-based tools for underserved populations.
\end{itemize}

\section{Selected Technical Writing}
\begin{itemize}
  \item \textit{BIO614-FinalProjectProposal} (Overleaf) --- formal manuscript on RNA secondary structure prediction with thermodynamic deep learning enhancements, reproducible evaluation, and biological interpretation.
  \item \textit{BIO550 RNA-seq Analysis} --- reproducible differential-expression pipeline with DESeq2, QC validation, and pathway-level biological interpretation.
  \item \textit{Equitable Bioinformatics: Enhancing Diagnostic Decision-Making through RNA and Biomarker Data} (ISTE-780 Project Report) --- fairness-aware ML framework applied to biological diagnostic algorithms.
  \item \textit{Big Data Medical Diagnosis: A Multi-Method Approach} (RIT MS CS graduate research) --- ensemble deep learning and ML methods for clinical classification on high-dimensional biomedical data.
\end{itemize}

\section{Awards \& Recognition}
\begin{itemize}
  \item NSF Robert Noyce Teacher Scholarship (2024--2026) --- full funding for M.S. in Teaching Computer Science K--12.
  \item MS International Scholarship (2012--2015) --- MESCYT \& RIT, awarded for academic excellence.
  \item Honorable Mention Technology Award (2012) --- NYS STEP/CSTEP, 2nd place capstone project.
  \item IEEE Alumni Member (2009--Present) --- Member No. 90613000.
\end{itemize}

\section{Technical Competencies}
\textbf{Programming:} Python, R, Java, C++, C\#, SQL, MATLAB, JavaScript, Bash\\
\textbf{Deep Learning \& ML:} PyTorch, TensorFlow, Keras; LSTM, CNNs, Transformer architectures; scikit-learn, XGBoost; uncertainty quantification; SHAP analysis\\
\textbf{Bioinformatics:} ViennaRNA, Nussinov algorithm, Turner parameters, DESeq2, RNA-seq pipelines, BLAST, differential expression analysis\\
\textbf{Data and Statistics:} pandas, NumPy, SciPy, statistical modeling, reproducible benchmarking\\
\textbf{Infrastructure:} Git/GitHub, Docker, Linux, GCP, AWS, LaTeX

\section{Selected Fit for KTH Deep Learning for Biological Systems Position}
\begin{itemize}
  \item Direct research experience applying deep learning and neural network methods to biological sequence and structure problems, including RNA structure prediction with thermodynamic enhancements (BIO614, Overleaf manuscript).
  \item RNA-seq differential-expression analysis background (BIO550) provides practical genomics pipeline experience relevant to downstream deep learning on transcriptomic data.
  \item Proven ability to produce formal, manuscript-quality research outputs in computational biology with reproducible implementations.
  \item Strong deep learning stack (PyTorch, TensorFlow, LSTM, CNN) with experience applying these architectures to biological and biomedical data beyond toy benchmarks.
  \item Industry experience (VIOME) building preprocessing and feature-extraction pipelines for microbiome sequencing data provides hands-on grounding in the challenges of real biological data at scale.
\end{itemize}

\end{document}
